\section{Całka Kurzweila--Stieltjesa}
\label{sec:kurzweil-stieltjes}
Całka Kurzweila--Stieltjesa (KS) łączy przyrosty funkcji całkującej z cechowaniami znanymi z całki Henstocka--Kurzweila. Jej przewaga nad całką Riemanna--Stieltjesa ujawnia się przede wszystkim przy wspólnych nieciągłościach funkcji podcałkowej i całkującej. W tej sekcji, oprócz kryteriów istnienia, wyprowadzimy wzory obliczeniowe i oszacowania błędu przybliżeń.

\subsection*{Definicja, jednoznaczność i podstawowe oszacowania}
\begin{definicja}[Całka Kurzweila--Stieltjesa]
Niech $a<b$ i niech $f,\alpha:[a,b]\to\R$ będą funkcjami o skończonych wartościach. \emph{Cechowaniem} nazywamy funkcję $\delta:[a,b]\to(0,\infty)$. Podział znaczony $\mathcal P=\{(\xi_i,[x_{i-1},x_i])\}_{i=1}^n$, gdzie $a=x_0<\cdots<x_n=b$ i $\xi_i\in[x_{i-1},x_i]$, jest \emph{$\delta$-drobny}, jeżeli $[x_{i-1},x_i]\subset(\xi_i-\delta(\xi_i),\xi_i+\delta(\xi_i))$ dla każdego $i$. Definiujemy
\[
 S_\alpha(f,\mathcal P):=\sum_{i=1}^n f(\xi_i)\bigl(\alpha(x_i)-\alpha(x_{i-1})\bigr).
\]
Liczba $I\in\R$ jest całką KS z $f$ względem $\alpha$, jeżeli dla każdego $\varepsilon>0$ istnieje cechowanie $\delta$ takie, że $|S_\alpha(f,\mathcal P)-I|<\varepsilon$ dla każdego $\delta$-drobnego podziału znaczonego $\mathcal P$. Piszemy wówczas $I=(\KS)\int_a^b f\,\mathrm d\alpha$.
\end{definicja}
Istnienie podziału $\delta$-drobnego dla każdego cechowania wynika z lematu Cousina (lemat 4.2). Gdy $\alpha(x)=x$, definicja jest dokładnie definicją całki Henstocka--Kurzweila. W szczególności nie wolno utożsamiać dowolnej całki KS z całką względem miary Stieltjesa: potrzebne są dodatkowe założenia.

\begin{twierdzenie}[Kryterium Cauchy'ego, liniowość i zgodność]
\label{tw:ks-cauchy}
Dla funkcji $f,\alpha:[a,b]\to\R$ o skończonych wartościach całka $(\KS)\int_a^b f\,\mathrm d\alpha$ istnieje wtedy i tylko wtedy, gdy dla każdego $\varepsilon>0$ istnieje cechowanie $\delta$ takie, że
$|S_\alpha(f,\mathcal P)-S_\alpha(f,\mathcal Q)|<\varepsilon$ dla dowolnych dwóch $\delta$-drobnych podziałów znaczonych $\mathcal P,\mathcal Q$. Całka jest jednoznaczna i liniowa osobno w $f$ i $\alpha$, o ile rozważane całki istnieją. Jeżeli istnieje właściwa całka Riemanna--Stieltjesa $\int_a^b f\,\mathrm d\alpha$, to istnieje również całka KS i ma tę samą wartość.
\end{twierdzenie}
\begin{proof}
Jeżeli sumy są bliskie liczbie $I$, ich różnica jest mniejsza od $2\varepsilon$. Odwrotnie, dla każdego $m\ge1$ wybieramy cechowanie $\delta_m$ zapewniające różnice sum mniejsze niż $2^{-m}$ i zastępujemy je przez $\min(\delta_1,\ldots,\delta_m)$. Wybieramy odpowiadające im podziały $\mathcal P_m$. Gdy $k\ge m$, obie sumy $S_\alpha(f,\mathcal P_k)$ i $S_\alpha(f,\mathcal P_m)$ są porównywalne przez cechowanie $\delta_m$, toteż tworzą ciąg Cauchy'ego. Niech $I$ będzie jego granicą. Dla dowolnego $\delta_m$-drobnego podziału $\mathcal P$ porównujemy jego sumę z $\mathcal P_k$ i przechodzimy z $k$ do nieskończoności; otrzymujemy $|S_\alpha(f,\mathcal P)-I|\le2^{-m}$. Jednoznaczność wynika z istnienia podziałów drobnych dla minimum dwóch cechowań. Liniowość dostajemy z liniowości sum i ponownego użycia minimum cechowań. Dla całki Riemanna--Stieltjesa wystarcza cechowanie stałe mniejsze niż dopuszczalna średnica podziału.
\end{proof}

\begin{definicja}[Wahanie całkowite i funkcja regulowana]
Dla $\alpha:[a,b]\to\R$ definiujemy $V_a^b(\alpha):=\sup_P\sum_i|\alpha(x_i)-\alpha(x_{i-1})|$, gdzie supremum przebiega wszystkie podziały $P$ przedziału $[a,b]$. Piszemy $\alpha\in\BV[a,b]$, jeśli $V_a^b(\alpha)<\infty$. Funkcja $f:[a,b]\to\R$ jest \emph{regulowana}, gdy istnieją skończone granice $f(t-)$ dla $t\in(a,b]$ i $f(t+)$ dla $t\in[a,b)$. Równoważnie $f$ jest jednostajną granicą funkcji schodkowych, przy czym wartości w skończonej liczbie węzłów można określać niezależnie od wartości na otwartych przedziałach.
\end{definicja}
\begin{proof}[Uzasadnienie równoważności]
Każda funkcja schodkowa ma granice jednostronne, a jednostajna granica funkcji posiadających takie granice spełnia jednostronny warunek Cauchy'ego. W przeciwnym kierunku ustalmy $\varepsilon>0$. W każdym punkcie $t$ mamy jednostronne otoczenia, na których oscylacja $f$ jest mniejsza niż $\varepsilon$. Lemat Cousina dostarcza skończony podział, którego wnętrza komórek można rozciąć w znacznikach, tak aby na każdym z otrzymanych otwartych przedziałów oscylacja $f$ była mniejsza niż $\varepsilon$. Na tych przedziałach podstawiamy stałe, a w węzłach zachowujemy dokładne wartości $f$. Powstała funkcja schodkowa przybliża $f$ jednostajnie z błędem co najwyżej $\varepsilon$.
\end{proof}

\begin{lemat}[Oszacowanie przez wahanie]
\label{lem:ks-var}
Jeżeli $\alpha\in\BV[a,b]$ i $f:[a,b]\to\R$ jest ograniczona, to dla każdego podziału znaczonego $\mathcal P$ zachodzi $|S_\alpha(f,\mathcal P)|\le\|f\|_\infty V_a^b(\alpha)$. Jeżeli całka KS istnieje, to
\[
 \left|(\KS)\int_a^b f\,\mathrm d\alpha\right|\le\|f\|_\infty V_a^b(\alpha).
\]
\end{lemat}
\begin{proof}
Stosujemy nierówność trójkąta do sumy $S_\alpha(f,\mathcal P)$, a następnie przechodzimy do granicy po podziałach $\delta$-drobnych.
\end{proof}

\subsection*{Skoki, wartości w punktach i funkcje schodkowe}
Dla regulowanej $g:[a,b]\to\R$ stosujemy oznaczenia $\Delta^-g(t)=g(t)-g(t-)$ dla $t\in(a,b]$, $\Delta^+g(t)=g(t+)-g(t)$ dla $t\in[a,b)$ i $\Delta g(t)=g(t+)-g(t-)$ dla $t\in(a,b)$. Na końcach przyjmujemy $\Delta^-g(a)=\Delta^+g(b)=0$ oraz $\Delta g(a)=\Delta^+g(a)$, $\Delta g(b)=\Delta^-g(b)$. Rozróżnienie skoku lewostronnego i prawostronnego jest istotne: $g(t)$ nie musi równać się żadnej granicy jednostronnej.

\begin{lemat}[Punktowy wzór wykonawczy]
\label{lem:ks-point}
Niech $g:[a,b]\to\R$ będzie regulowana i $c\in[a,b]$. Dla dowolnej liczby $q\in\R$ zachodzi
\[
 (\KS)\int_a^b q\ind{\{c\}}(t)\,\mathrm dg(t)=q\Delta g(c).
\]
W szczególności wartość funkcji podcałkowej w punkcie skoku nie może być pomijana.
\end{lemat}
\begin{proof}
Rozpatrzmy $a<c<b$. Dla $t\ne c$ wybieramy $\delta(t)<|t-c|/2$; wartość $\delta(c)>0$ dobieramy tak małą, aby $|g(u)-g(c-)|<\varepsilon$ dla $c-\delta(c)<u<c$ i $|g(v)-g(c+)|<\varepsilon$ dla $c<v<c+\delta(c)$. W podziale drobnym każda komórka, której wnętrze zawiera $c$, musi mieć znacznik $c$. Jeżeli $c$ jest węzłem podziału, jedynie dwie przyległe do niego komórki mogą mieć znacznik $c$; obie rzeczywiście muszą go mieć przy dodatniej długości. Niezerowe składniki sumy teleskopują więc do $q(g(v)-g(u))$, gdzie $u\le c\le v$, a odległości odpowiednich końców od $c$ są mniejsze niż $\delta(c)$. Wartość sumy różni się od $q(g(c+)-g(c-))$ co najwyżej o $2|q|\varepsilon$. Gdy $c=a$ lub $c=b$, pozostaje tylko jedna z dwóch stron, co daje odpowiednio $q\Delta^+g(a)$ i $q\Delta^-g(b)$. W razie $q=0$ twierdzenie jest natychmiastowe.
\end{proof}

\begin{lemat}[Wskaźniki przedziałów]
\label{lem:ks-indicators}
Jeśli $g:[a,b]\to\R$ jest regulowana, to dla $a\le u<v\le b$ istnieją całki KS z funkcji charakterystycznych $[u,v]$, $(u,v)$, $[u,v)$ i $(u,v]$. Przykładowo dla $a<u<v<b$ otrzymujemy
\[
 \begin{aligned}
 (\KS)\int_a^b\ind{(u,v)}\,\mathrm dg&=g(v-)-g(u+),\\
 (\KS)\int_a^b\ind{[u,v]}\,\mathrm dg&=g(v+)-g(u-),\\
 (\KS)\int_a^b\ind{[u,v)}\,\mathrm dg&=g(v-)-g(u-),\\
 (\KS)\int_a^b\ind{(u,v]}\,\mathrm dg&=g(v+)-g(u+).
 \end{aligned}
\]
Gdy $u=a$ lub $v=b$, symbole $g(a-)$ i $g(b+)$ we wzorach zastępujemy odpowiednio $g(a)$ i $g(b)$.
\end{lemat}
\begin{proof}
Dla funkcji $\ind{(u,v)}$ wybieramy cechowanie tak, aby przedziały o znacznikach różnych od $u,v$ nie przecinały tych punktów. Komórki mające znaczniki $u$ lub $v$ wnoszą zero, a wszystkie komórki ze znacznikami wewnątrz $(u,v)$ tworzą jeden blok. Ich przyrosty teleskopują do $g(r)-g(s)$, gdzie $s\to u+$ i $r\to v-$, skąd pierwszy wzór. Następnie korzystamy z tożsamości funkcji charakterystycznych, np. $\ind{[u,v]}=\ind{(u,v)}+\ind{\{u\}}+\ind{\{v\}}$, oraz z lematu~\ref{lem:ks-point}. Pozostałe wzory uzyskujemy analogicznie.
\end{proof}

\begin{twierdzenie}[Całkowanie względem funkcji schodkowej]
\label{tw:ks-step}
Niech $g:[a,b]\to\R$ będzie funkcją schodkową, mającą skończony zbiór punktów nieciągłości $D$, i niech $f:[a,b]\to\R$ będzie dowolną funkcją o skończonych wartościach. Wtedy całka KS istnieje i
\[
 (\KS)\int_a^b f\,\mathrm dg
 =f(a)\Delta^+g(a)+\sum_{c\in D\cap(a,b)}f(c)\Delta g(c)+f(b)\Delta^-g(b).
\]
\end{twierdzenie}
\begin{proof}
Zapisujemy $g$ jako sumę stałej oraz skończonej liczby funkcji $\ind{[c,b]}$ i $\ind{(c,b]}$, uwzględniając oddzielnie skoki przy końcach. Z lematu~\ref{lem:ks-point} wynika dla każdej z tych funkcji całkujących tożsamość $\int f\,\mathrm d\ind{[c,b]}=f(c)$ i $\int f\,\mathrm d\ind{(c,b]}=f(c)$. Liniowość względem integratora daje wskazany wzór. Nie potrzebujemy ograniczoności ani regularności $f$, ponieważ występuje skończenie wiele wartości $f(c)$.
\end{proof}

\begin{przyklad}[Wspólny skok i nieistnienie zwykłej całki Stieltjesa]
Niech $a<c<b$ i $H=\ind{[c,b]}$. Wtedy dla każdej funkcji $f$ o skończonych wartościach zachodzi $(\KS)\int_a^b f\,\mathrm dH=f(c)$. Stąd $(\KS)\int_a^b H\,\mathrm dH=1$, a $(\KS)\int_a^b\ind{(c,b]}\,\mathrm dH=0$. Pierwsza z tych całek nie istnieje jako zwykła całka Riemanna--Stieltjesa, ponieważ swobodny wybór znacznika na komórce przecinającej $c$ daje różne sumy. Całka KS rozstrzyga tę niejednoznaczność za pomocą cechowania.
\end{przyklad}

\begin{przyklad}[Punktowy impuls integratora]
Niech $g=\ind{\{c\}}$ dla $a<c<b$. Wtedy $\Delta^-g(c)=1$ i $\Delta^+g(c)=-1$, zatem $\Delta g(c)=0$. Dla każdego regulowanego $f$ otrzymujemy $(\KS)\int_a^b f\,\mathrm dg=0$. Oznacza to, że dwie przeciwne zmiany funkcji całkującej, następujące po obu stronach tego samego punktu, znoszą się. Wartość $g(c)$ może więc nie wpływać na całkę, mimo że oddzielne skoki jednostronne są niezerowe.
\end{przyklad}

\subsection*{Dwa podstawowe twierdzenia o istnieniu}
\begin{twierdzenie}[Funkcja regulowana i całkująca o wahaniu ograniczonym]
\label{tw:ks-reg-bv}
Jeśli $f:[a,b]\to\R$ jest regulowana, a $g\in\BV[a,b]$, to całka $(\KS)\int_a^b f\,\mathrm dg$ istnieje. Ponadto
$|(\KS)\int_a^b f\,\mathrm dg|\le\|f\|_\infty V_a^b(g)$.
\end{twierdzenie}
\begin{proof}
Wybieramy funkcje schodkowe $f_n$ z $\|f_n-f\|_\infty\to0$. Każda całka $I_n=(\KS)\int_a^b f_n\,\mathrm dg$ istnieje na mocy lematu~\ref{lem:ks-indicators}, ponieważ $f_n$ jest skończoną kombinacją wskaźników przedziałów i punktów. Z lematu~\ref{lem:ks-var} i liniowości wynika $|I_n-I_m|\le\|f_n-f_m\|_\infty V_a^b(g)$. Niech $I=\lim_n I_n$. Dla dowolnego podziału znaczonego $\mathcal P$ mamy
\[
 |S_g(f,\mathcal P)-I|\le
 \|f-f_n\|_\infty V_a^b(g)+|S_g(f_n,\mathcal P)-I_n|+|I_n-I|.
\]
Dla ustalonego $\varepsilon$ najpierw wybieramy $n$ zapewniające małość pierwszego i trzeciego składnika, a następnie cechowanie dla $f_n$. Ostatnia nierówność wynika z lematu~\ref{lem:ks-var} po przejściu do granicy.
\end{proof}

\begin{lemat}[Nierówność sumowania przez części]
\label{lem:ks-summation}
Niech $f\in\BV[a,b]$ i niech $g:[a,b]\to\R$ będzie ograniczona. Dla każdego podziału znaczonego $\mathcal P$ zachodzi
\[
 |S_g(f,\mathcal P)|\le
 \bigl(2\|f\|_\infty+V_a^b(f)\bigr)\|g\|_\infty.
\]
\end{lemat}
\begin{proof}
Oznaczmy $f_i=f(\xi_i)$, $g_i=g(x_i)$. Przekształcenie skończonej sumy daje
$S_g(f,\mathcal P)=f_ng_n-f_1g_0+\sum_{i=1}^{n-1}(f_i-f_{i+1})g_i$. Ponieważ znaczniki są uporządkowane, $\xi_i\le\xi_{i+1}$, mamy $\sum_{i=1}^{n-1}|f_i-f_{i+1}|\le V_a^b(f)$. Oszacowanie wynika z nierówności trójkąta.
\end{proof}

\begin{twierdzenie}[Całkująca regulowana i funkcja podcałkowa o wahaniu ograniczonym]
\label{tw:ks-bv-reg}
Jeśli $f\in\BV[a,b]$ i $g:[a,b]\to\R$ jest regulowana, to całka $(\KS)\int_a^b f\,\mathrm dg$ istnieje. Jeżeli $g,h$ są regulowane, to
\[
 \left|(\KS)\int_a^b f\,\mathrm dg-(\KS)\int_a^b f\,\mathrm dh\right|
 \le\bigl(2\|f\|_\infty+V_a^b(f)\bigr)\|g-h\|_\infty.
\]
\end{twierdzenie}
\begin{proof}
Funkcję $g$ przybliżamy jednostajnie funkcjami schodkowymi $g_n$. Całki $I_n=(\KS)\int_a^b f\,\mathrm dg_n$ istnieją na mocy twierdzenia~\ref{tw:ks-step}. Z lematu~\ref{lem:ks-summation}, zastosowanego do $g_n-g_m$, otrzymujemy $|I_n-I_m|\le K_f\|g_n-g_m\|_\infty$, gdzie $K_f=2\|f\|_\infty+V_a^b(f)$. Istnieje zatem $I=\lim_nI_n$. Dla dowolnego podziału $\mathcal P$ mamy
$|S_g(f,\mathcal P)-I|\le K_f\|g-g_n\|_\infty+|S_{g_n}(f,\mathcal P)-I_n|+|I_n-I|$. Wybór najpierw $n$, a następnie cechowania dla $g_n$, dowodzi istnienia całki. To samo oszacowanie dla całki różnicy $g-h$ daje drugą tezę.
\end{proof}

\begin{uwaga}[Właściwa dziedzina obu twierdzeń]
Twierdzenia~\ref{tw:ks-reg-bv} i~\ref{tw:ks-bv-reg} są różne. Pierwsze pozwala na funkcję podcałkową regulowaną o nieskończonym wahaniu, o ile całkująca ma wahanie ograniczone; drugie pozwala na całkującą regulowaną o nieskończonym wahaniu, gdy funkcja podcałkowa ma wahanie ograniczone. Żadne z tych twierdzeń nie wymaga ciągłości obu funkcji ani braku wspólnych skoków.
\end{uwaga}

\begin{twierdzenie}[Konstruktywny rozkład funkcji o wahaniu ograniczonym]
\label{tw:ks-bv-split}
Dla $g\in\BV[a,b]$ zdefiniujmy
\[
 g_{\mathrm{sk}}(t)=
 \sum_{c\in(a,b]}\Delta^-g(c)\ind{[c,b]}(t)
 +\sum_{c\in[a,b)}\Delta^+g(c)\ind{(c,b]}(t).
\]
Oba szeregi są bezwzględnie i jednostajnie zbieżne, a funkcja $g_{\mathrm c}:=g-g(a)-g_{\mathrm{sk}}$ jest ciągła i ma wahanie ograniczone. Dla każdej regulowanej $f:[a,b]\to\R$ zachodzi
\[
 \begin{aligned}
 (\KS)\int_a^b f\,\mathrm dg
 ={}&(\KS)\int_a^b f\,\mathrm dg_{\mathrm c}
 +f(a)\Delta^+g(a)\\
 &+\sum_{c\in(a,b)}f(c)\bigl(\Delta^-g(c)+\Delta^+g(c)\bigr)
 +f(b)\Delta^-g(b).
 \end{aligned}
\]
W szczególności dla funkcji całkującej złożonej wyłącznie ze skoków całka jest szeregiem wartości $f$ w punktach skoków, niezależnie od ich liczby.
\end{twierdzenie}
\begin{proof}
Dla każdego skończonego zbioru punktów $E\subset[a,b]$ suma modułów wszystkich skoków jednostronnych $g$ w tych punktach nie przekracza $V_a^b(g)$. Aby to zobaczyć, wybieramy parami rozłączne, odpowiednio małe otoczenia punktów $E$, a następnie podziały uwzględniające punkty $c$ i po jednym punkcie z każdej dozwolonej strony $c$. W definicji wahania przechodzimy do granic jednostronnych. Biorąc supremum po skończonych $E$, otrzymujemy
$\sum_{c\in(a,b]}|\Delta^-g(c)|+\sum_{c\in[a,b)}|\Delta^+g(c)|\le V_a^b(g)$. Zatem szeregi definiujące $g_{\mathrm{sk}}$ są absolutnie i jednostajnie zbieżne.

Skok lewostronny funkcji $g_{\mathrm{sk}}$ w $c\in(a,b]$ jest równy $\Delta^-g(c)$, a skok prawostronny w $c\in[a,b)$ wynosi $\Delta^+g(c)$. Ponieważ $g$ jest regulowana, różnica $g_{\mathrm c}$ nie ma żadnych skoków jednostronnych i dlatego jest ciągła. Ma również wahanie ograniczone, jako różnica funkcji o wahaniu ograniczonym. Dla skończonych sum wzór na całkę wynika z twierdzenia~\ref{tw:ks-step}. Przy zwiększaniu liczby składników ogon ma wahanie ograniczone przez sumę modułów pominiętych skoków, a więc jego całka ma moduł nie większy niż $\|f\|_\infty$ razy ten ogon, na mocy lematu~\ref{lem:ks-var}. Przejście do granicy dowodzi wzoru.
\end{proof}

\subsection*{Wzory przez miarę Stieltjesa i rozkład Lebesgue'a}
\begin{twierdzenie}[Miara kanonicznie związana z funkcją $\BV$]
\label{tw:ks-measure}
Niech $g\in\BV[a,b]$. Definiujemy $G:[a,b]\to\R$ przez $G(t)=g(t+)-g(a)$ dla $a\le t<b$ oraz $G(b)=g(b)-g(a)$. Istnieje dokładnie jedna skończona \emph{podpisana miara borelowska} $\mu_g$ na $[a,b]$ spełniająca $\mu_g([a,t])=G(t)$ dla każdego $t\in[a,b]$. Jej atomy są dane wzorami
\[
 \mu_g(\{a\})=\Delta^+g(a),\qquad
 \mu_g(\{c\})=\Delta g(c)\ \ (a<c<b),\qquad
 \mu_g(\{b\})=\Delta^-g(b).
\]
Jeśli $f$ jest regulowana, to $f$ jest całkowalna względem $|\mu_g|$ i
\[
 (\KS)\int_a^b f\,\mathrm dg=\int_{[a,b]}f(t)\,\mathrm d\mu_g(t),
 \qquad |\mu_g|([a,b])\le V_a^b(g).
\]
\end{twierdzenie}
\begin{proof}
Funkcja $G$ jest prawostronnie ciągła na $[a,b)$ i ma wahanie ograniczone. Rozkłada się na różnicę dwóch funkcji niemalejących i prawostronnie ciągłych (stosujemy funkcję wahania całkowitego i rozkład Jordana). Każda z tych funkcji wyznacza skończoną miarę Stieltjesa; ich różnica daje $\mu_g$. Jednoznaczność wynika z jednoznaczności miar na półpierścieniu przedziałów $(s,t]$ wraz z jednoelementowym zbiorem $\{a\}$. Wzory na atomy otrzymujemy przez różnice granic jednostronnych $G$. Nierówność dla wahania miary wynika z oszacowania wszystkich sum bezwzględnych przyrostów przez $V_a^b(g)$.

Dla funkcji $\ind{\{c\}}$ równość całek jest lematem~\ref{lem:ks-point}; dla wskaźników czterech typów przedziałów otrzymujemy ją z lematu~\ref{lem:ks-indicators} oraz wzorów na masy przedziałów. Z liniowości wynika równość dla każdej funkcji schodkowej. Wybieramy funkcje schodkowe $f_n\to f$ jednostajnie. Z twierdzenia~\ref{tw:ks-reg-bv} i lematu~\ref{lem:ks-var} całki KS z $f_n$ zbiegają do całki KS z $f$, zaś całki względem $\mu_g$ zbiegają dzięki nierówności $|\int(f_n-f)\,\mathrm d\mu_g|\le\|f_n-f\|_\infty|\mu_g|([a,b])$. Stąd teza.
\end{proof}

\begin{uwaga}[Konwencja na końcach i równość prawie wszędzie]
Jeśli $g$ jest niemalejąca i prawostronnie ciągła na $[a,b]$, to $\mu_g(\{a\})=0$ i $\mu_g((s,t])=g(t)-g(s)$; jest to konwencja używana wcześniej przy całce Lebesgue'a--Stieltjesa. Dla funkcji, która nie jest prawostronnie ciągła, nie wolno stosować tego samego wzoru bez uwzględnienia przesunięcia do $g(t+)$. Równość $f=h$ prawie wszędzie względem miary Lebesgue'a nie wystarcza do równości całek KS względem $g$ z atomami: trzeba uwzględnić wartości w punktach, którym $\mu_g$ przypisuje masę.
\end{uwaga}

\begin{twierdzenie}[Wzór: część absolutnie ciągła, osobliwa i skoki]
\label{tw:ks-decomp}
Niech $g\in\BV[a,b]$ i $f:[a,b]\to\R$ będzie regulowana. Niech rozkład Lebesgue'a podpisanej miary $\mu_g$ względem $\lambda$ ma postać
$\mu_g=q\lambda+\mu_{\mathrm{sc}}+\mu_{\mathrm{at}}$, gdzie $q\in L^1([a,b],\lambda)$, $\mu_{\mathrm{sc}}$ jest bezatomową miarą podpisaną osobliwą względem $\lambda$, a $\mu_{\mathrm{at}}$ jest częścią czysto atomową. Wtedy
\[
 \begin{aligned}
 (\KS)\int_a^b f\,\mathrm dg
 ={}&\int_a^b f(t)q(t)\,\mathrm d\lambda(t)
 +\int_{[a,b]}f\,\mathrm d\mu_{\mathrm{sc}}\\
 &+f(a)\Delta^+g(a)+\sum_{c\in(a,b)}f(c)\Delta g(c)
 +f(b)\Delta^-g(b).
 \end{aligned}
\]
Szereg jest bezwzględnie zbieżny. Jeżeli $g$ nie ma osobliwej części ciągłej, drugi składnik znika; jeśli $g$ jest absolutnie ciągła, znikają również wszystkie składniki skokowe i $q=g'$ prawie wszędzie.
\end{twierdzenie}
\begin{proof}
Z twierdzenia~\ref{tw:ks-measure} całka KS jest całką względem skończonej miary podpisanej $\mu_g$. Stosujemy rozkład Lebesgue'a jej części bezatomowej względem $\lambda$ i liniowość całki. Każda skończona miara podpisana ma co najwyżej przeliczalnie wiele atomów niezerowej masy, a suma modułów ich mas nie przekracza $|\mu_g|([a,b])\le V_a^b(g)$. Stąd $\sum|f(c)\Delta g(c)|\le\|f\|_\infty V_a^b(g)$ i bezwzględna zbieżność. Ostatnia część wynika ze wzoru Newtona--Leibniza dla funkcji absolutnie ciągłej $g(t)=g(a)+\int_a^t g'(s)\,\mathrm d\lambda(s)$.
\end{proof}

\begin{wniosek}[Wzór dla funkcji absolutnie ciągłej z impulsami]
\label{wn:ks-ac-jumps}
Niech $h\in\AC[a,b]$, $a<c_1<\cdots<c_m<b$, $A_j,B_j\in\R$ oraz
$g(t)=h(t)+\sum_{j=1}^m\bigl(A_j\ind{[c_j,b]}(t)+B_j\ind{(c_j,b]}(t)\bigr)$. Dla każdej regulowanej $f$ zachodzi
\[
 (\KS)\int_a^b f\,\mathrm dg
 =\int_a^b f(t)h'(t)\,\mathrm d\lambda(t)
 +\sum_{j=1}^m(A_j+B_j)f(c_j).
\]
Gdy $h\in C^1[a,b]$, pierwsza całka jest zwykłą całką Riemanna z funkcji regulowanej $fh'$.
\end{wniosek}
\begin{proof}
Dla części $h$ stosujemy twierdzenie~\ref{tw:ks-decomp}; dla każdej funkcji skokowej stosujemy twierdzenie~\ref{tw:ks-step}. Liniowość względem funkcji całkującej kończy dowód.
\end{proof}

\begin{przyklad}[Skończona liczba impulsów: pełny rachunek]
Na $[0,1]$ niech $g(t)=t+3\ind{[1/3,1]}(t)-2\ind{(2/3,1]}(t)$ oraz $f(t)=t^2$. Ponieważ funkcja $t\mapsto t$ jest absolutnie ciągła, a skoki $g$ mają wielkości odpowiednio $3$ i $-2$, otrzymujemy
\[
 (\KS)\int_0^1 t^2\,\mathrm dg(t)
 =\int_0^1 t^2\,\mathrm d\lambda(t)
 +3\left(\frac13\right)^2-2\left(\frac23\right)^2
 =\frac13+\frac13-\frac89=-\frac29.
\]
Wartość $g(2/3)$ jest tutaj inna niż $g((2/3)+)$, ale wzór na całkę korzysta z całkowitego skoku $\Delta g(2/3)$.
\end{przyklad}

\begin{przyklad}[Dwa atomy na końcach przedziału]
Na $[0,1]$ niech $g(t)=\ind{(0,1]}(t)+2\ind{\{1\}}(t)$ i $f(t)=1+t$. Funkcja $g$ ma atom o masie $1$ w punkcie $0$ i atom o masie $2$ w punkcie $1$, mimo że wewnątrz przedziału jest stała. Z lematu~\ref{lem:ks-point} otrzymujemy
$(\KS)\int_0^1 f\,\mathrm dg=f(0)\cdot1+f(1)\cdot2=5$.
Ten rachunek pokazuje, dlaczego w twierdzeniach trzeba osobno określić sposób traktowania końców przedziału.
\end{przyklad}

\begin{przyklad}[Nieskończenie wiele impulsów i oszacowanie ogona]
Niech $c_n=1-2^{-n}$, $g(t)=\sum_{n\ge1}2^{-n}\ind{[c_n,1]}(t)$ na $[0,1]$ oraz $f(t)=t$. Wówczas $g$ jest niemalejąca, $V_0^1(g)=1$ i
\[
 (\KS)\int_0^1 t\,\mathrm dg(t)
 =\sum_{n=1}^\infty2^{-n}(1-2^{-n})
 =1-\frac13=\frac23.
\]
Jeśli liczymy tylko pierwszych $N$ impulsów, błąd nie przekracza $\|f\|_\infty\sum_{n>N}2^{-n}=2^{-N}$. Jest to konkretny, certyfikowany algorytm przybliżania całki.
\end{przyklad}

\begin{przyklad}[Miara Cantora: brak wzoru $g'\,\mathrm dx$]
Niech $C:[0,1]\to[0,1]$ będzie funkcją Cantora. Funkcja ta jest ciągła, niemalejąca, $C'=0$ prawie wszędzie i $C(1)-C(0)=1$, więc jej miara Stieltjesa jest osobliwa ciągła. Dla $f(t)=t$ dostajemy $(\KS)\int_0^1 t\,\mathrm dC(t)=1/2$ (z symetrii miary Cantora), a dla $f(t)=t^2$ wartość $3/8$ (obliczenie momentów w rozdziale 10). Zastąpienie $\mathrm dC$ przez $C'(t)\,\mathrm d\lambda(t)$ dałoby błędnie zero. Dlatego w twierdzeniu~\ref{tw:ks-decomp} osobliwa część ciągła jest niezbędna.
\end{przyklad}

\subsection*{Całkowanie przez części ze wszystkimi poprawkami skokowymi}
\begin{twierdzenie}[Wzór całkowania przez części]
\label{tw:ks-ibp}
Niech $f,g:[a,b]\to\R$ będą regulowane i niech przynajmniej jedna z funkcji $f,g$ ma wahanie ograniczone. Wtedy obie całki KS istnieją i
\[
 \begin{aligned}
 (\KS)\int_a^b f\,\mathrm dg+(\KS)\int_a^b g\,\mathrm df
 ={}&f(b)g(b)-f(a)g(a)\\
 &+\sum_{c\in[a,b]}\bigl(\Delta^-f(c)\Delta^-g(c)
 -\Delta^+f(c)\Delta^+g(c)\bigr).
 \end{aligned}
\]
Przyjmujemy $\Delta^-f(a)=\Delta^-g(a)=0$ i $\Delta^+f(b)=\Delta^+g(b)=0$. Szereg korekcyjny jest bezwzględnie zbieżny. W szczególności, jeśli jedna z funkcji jest ciągła, cały wyraz korekcyjny znika.
\end{twierdzenie}
\begin{proof}
\emph{Krok 1: obliczenie przy jednym wspólnym skoku.} Ustalmy $c\in(a,b)$ i zapiszmy $f_-=f(c-)$, $f_0=f(c)$, $f_+=f(c+)$ oraz analogicznie $g_-,g_0,g_+$. Punktowy wkład do sumy dwóch całek wynosi $f_0(g_+-g_-)+g_0(f_+-f_-)$. Tymczasem bezpośrednie wymnożenie daje
\[
 f_+g_+-f_-g_-=f_0(g_+-g_-)+g_0(f_+-f_-)
 +(f_+-f_0)(g_+-g_0)-(f_0-f_-)(g_0-g_-).
\]
Po przeniesieniu dwóch ostatnich składników na drugą stronę otrzymujemy wyraz $\Delta^-f(c)\Delta^-g(c)-\Delta^+f(c)\Delta^+g(c)$. Na końcach $a,b$ występuje tylko jedna z dwóch połówek skoku.

\emph{Krok 2: funkcje schodkowe.} Wybieramy wspólny skończony zbiór węzłów obu funkcji. Na otwartych przedziałach pomiędzy nimi obie funkcje są stałe. Całki liczymy z twierdzenia~\ref{tw:ks-step}, sumując wkłady w węzłach. Powyższa tożsamość algebraiczna po zsumowaniu teleskopuje się do różnicy $f(b)g(b)-f(a)g(a)$ i sumy poprawek.

\emph{Krok 3: przejście do funkcji regulowanych.} Możemy bez straty ogólności założyć $f\in\BV[a,b]$. Funkcję $g$ przybliżamy jednostajnie funkcjami schodkowymi $g_n$. Dla ustalonej funkcji schodkowej $g_n$ otrzymujemy tożsamość najpierw dla schodkowych przybliżeń $f_m$ funkcji $f$ i przechodzimy do granicy przy $m\to\infty$: całka $\int f_m\,\mathrm dg_n$ jest ciągła w normie jednostajnej z oszacowaniem przez $V(g_n)$, a całka $\int g_n\,\mathrm df_m$ jest ciągła względem normy jednostajnej integratora dzięki lematowi~\ref{lem:ks-summation}, zastosowanemu do $g_n$ jako funkcji podcałkowej o wahaniu ograniczonym. Sumy poprawek zawierają tylko skończenie wiele punktów skokowych $g_n$, więc również przechodzą do granicy. Istnienie potrzebnych schodkowych przybliżeń $f_m\to f$ wynika z regulowania $f$.

Przy $n\to\infty$ całki $\int f\,\mathrm dg_n$ zbiegają do $\int f\,\mathrm dg$ z twierdzenia~\ref{tw:ks-bv-reg}; całki $\int g_n\,\mathrm df$ zbiegają dzięki lematowi~\ref{lem:ks-var}. Skoki $g_n$ zbiegają jednostajnie do odpowiednich jednostronnych skoków $g$, a sumę ich iloczynów ze skokami $f$ kontrolujemy przez $V(f)\|g_n-g\|_\infty$: przykładowo $\sum_c|\Delta^-f(c)\Delta^-(g_n-g)(c)|\le2\|g_n-g\|_\infty V(f)$. Druga suma podlega temu samemu oszacowaniu. Wyrazy korekcyjne zbiegają więc bezwzględnie. Dowód zakończony.
\end{proof}

\begin{wniosek}[Klasyczna postać i wzór dla funkcji absolutnie ciągłej]
\label{wn:ks-ibp-ac}
Jeśli $f\in\AC[a,b]$, a $g:[a,b]\to\R$ jest regulowana, to
\[
 (\KS)\int_a^b f\,\mathrm dg
 =f(b)g(b)-f(a)g(a)-\int_a^b g(t)f'(t)\,\mathrm d\lambda(t).
\]
Jeśli $f\in C^1[a,b]$, ostatnią całkę można rozumieć jako zwykłą całkę Riemanna. Nie wymagamy w tym wniosku, aby $g$ miała wahanie ograniczone.
\end{wniosek}
\begin{proof}
Funkcja absolutnie ciągła ma wahanie ograniczone i jest ciągła, więc wzór z twierdzenia~\ref{tw:ks-ibp} nie ma poprawek skokowych. Ponadto $\mathrm df=f'(t)\,\mathrm d\lambda(t)$; z twierdzenia~\ref{tw:ks-decomp}, zastosowanego do całki $\int g\,\mathrm df$ (integrator $f\in\BV$, funkcja podcałkowa $g$ jest regulowana), otrzymujemy $\int g\,\mathrm df=\int gf'\,\mathrm d\lambda$. Iloczyn $gf'$ jest całkowalny w sensie Lebesgue'a, ponieważ $g$ jest ograniczona, a $f'\in L^1$.
\end{proof}

\begin{przyklad}[Rachunek bez znajomości pochodnej funkcji całkującej]
Niech $g:[0,1]\to\R$ będzie dowolną funkcją regulowaną (nie zakładamy $g\in\BV$). Dla $f(t)=t^2$ otrzymujemy
\[
 (\KS)\int_0^1 t^2\,\mathrm dg(t)=g(1)-2\int_0^1 t\,g(t)\,\mathrm d\lambda(t).
\]
Ponieważ funkcja regulowana jest ograniczona i Riemanna całkowalna, ostatnia całka jest zwykłą całką Riemanna. W szczególności wzór działa dla funkcji $g$ oscylującej z nieskończonym wahaniem, o ile ma granice jednostronne.
\end{przyklad}

\begin{przyklad}[Całkująca o nieskończonym wahaniu]
Dla $t\in(0,1]$ niech $g(t)=t\sin(1/t^2)$, a $g(0)=0$. Funkcja $g$ jest ciągła, ale nie ma wahania ograniczonego: na punktach $t_n=(\pi(n+1/2))^{-1/2}$ przyjmuje naprzemienne wartości $\pm t_n$, a suma sąsiednich oscylacji jest rozbieżna. Pomimo tego dla $f(t)=t$ całka KS istnieje na mocy twierdzenia~\ref{tw:ks-bv-reg}, a wzór całkowania przez części daje
\[
 (\KS)\int_0^1 t\,\mathrm dg(t)
 =\sin1-\int_0^1 t\sin(1/t^2)\,\mathrm d\lambda(t)
 =\sin1-\frac12\int_1^\infty\frac{\sin u}{u^2}\,\mathrm du.
\]
Ostatnia całka jest bezwzględnie zbieżna i może być obliczana metodami numerycznymi. Nie potrzebujemy ani wahania ograniczonego funkcji $g$, ani klasycznej całki Riemanna--Stieltjesa.
\end{przyklad}

\begin{przyklad}[Dlaczego potrzebujemy wyrazu korekcyjnego]
Dla $H=\ind{[c,b]}$ mamy $\Delta^-H(c)=1$ i $\Delta^+H(c)=0$. Twierdzenie~\ref{tw:ks-ibp} z $f=g=H$ daje
$2(\KS)\int_a^b H\,\mathrm dH=H(b)^2-H(a)^2+1=2$. Dla $K=\ind{(c,b]}$ mamy natomiast $\Delta^-K(c)=0$, $\Delta^+K(c)=1$, zatem $2(\KS)\int_a^b K\,\mathrm dK=1-1=0$. Oba wyniki są zgodne z bezpośrednim wzorem $\int f\,\mathrm dH=f(c)$ i $\int f\,\mathrm dK=f(c)$.
\end{przyklad}

\subsection*{Całka nieoznaczona, podstawienie i błędy numeryczne}
\begin{twierdzenie}[Skoki funkcji pierwotnej KS]
\label{tw:ks-primitive}
Niech $f$ będzie regulowana i $g\in\BV[a,b]$. Zdefiniujmy $F(x)=(\KS)\int_a^x f(t)\,\mathrm dg(t)$ dla $x\in[a,b]$. Wtedy $F\in\BV[a,b]$, $F(a)=0$, $V_a^b(F)\le\|f\|_\infty V_a^b(g)$, a ponadto
\[
 F(c+)-F(c)=f(c)\Delta^+g(c),\qquad
 F(c)-F(c-)=f(c)\Delta^-g(c)
\]
odpowiednio dla $c\in[a,b)$ i $c\in(a,b]$.
\end{twierdzenie}
\begin{proof}
Z addytywności całki po przyległych przedziałach i z lematu~\ref{lem:ks-var} mamy dla każdego podziału $a=x_0<\cdots<x_n=b$ nierówność $\sum_i|F(x_i)-F(x_{i-1})|\le\|f\|_\infty\sum_i V_{x_{i-1}}^{x_i}(g)=\|f\|_\infty V_a^b(g)$. Dla $x\downarrow c$ rozdzielamy całkę $\int_c^x f\,\mathrm dg$ na atom w punkcie $c$, mający wartość $f(c)(g(c+)-g(c))$, oraz całkę względem pozostałej miary na $(c,x]$. Jej moduł jest ograniczony przez $\|f\|_\infty|\mu_g|((c,x])\to0$. Otrzymujemy pierwszy wzór. Dowód wzoru lewostronnego jest analogiczny.
\end{proof}

\begin{twierdzenie}[Podstawienie w przedziale]
\label{tw:ks-substitution}
Niech $\varphi:[A,B]\to[a,b]$ będzie ściśle rosnącym homeomorfizmem, $\varphi(A)=a$, $\varphi(B)=b$. Niech $f,g:[a,b]\to\R$ będą funkcjami o skończonych wartościach. Wtedy całka KS po jednej stronie istnieje wtedy i tylko wtedy, gdy istnieje po drugiej, oraz
\[
 (\KS)\int_a^b f(t)\,\mathrm dg(t)
 =(\KS)\int_A^B f(\varphi(s))\,\mathrm d(g\circ\varphi)(s).
\]
\end{twierdzenie}
\begin{proof}
Każdy podział znaczony $\mathcal Q=\{(\eta_i,[s_{i-1},s_i])\}$ przedziału $[A,B]$ przenosi się na podział $\mathcal P=\{(\varphi(\eta_i),[\varphi(s_{i-1}),\varphi(s_i)])\}$ przedziału $[a,b]$. Obie sumy KS są równe składnik po składniku. Niech $\delta$ będzie cechowaniem na $[a,b]$. Dzięki ciągłości $\varphi$ w każdym punkcie $s$ można wybrać $\widetilde\delta(s)>0$ tak, aby $|u-s|<\widetilde\delta(s)$ implikowało $|\varphi(u)-\varphi(s)|<\delta(\varphi(s))$. Wtedy każdy podział $\widetilde\delta$-drobny przenosi się na podział $\delta$-drobny, zatem istnienie i wartość całki po lewej implikują istnienie i tę samą wartość po prawej. Dla implikacji odwrotnej powtarzamy argument z $\varphi^{-1}$.
\end{proof}

\begin{twierdzenie}[Podstawienie całki nieoznaczonej]
\label{tw:ks-chain}
Niech $f$ będzie regulowana, $g\in\BV[a,b]$, $F(x)=(\KS)\int_a^x f\,\mathrm dg$ i niech $h:[a,b]\to\R$ będzie regulowana. Wówczas
\[
 (\KS)\int_a^b h(t)\,\mathrm dF(t)
 =(\KS)\int_a^b h(t)f(t)\,\mathrm dg(t).
\]
\end{twierdzenie}
\begin{proof}
Na mocy twierdzenia~\ref{tw:ks-measure} funkcja $F$ ma wahanie ograniczone, a odpowiadająca jej miara Stieltjesa spełnia $\mu_F(E)=\int_E f\,\mathrm d\mu_g$ dla każdego zbioru borelowskiego $E\subset[a,b]$. Wystarczy zweryfikować tę tożsamość na przedziałach $[a,x]$, uwzględniając prawostronne wartości $F(x+)$ i $g(x+)$; wówczas wynika ona z definicji $F$ oraz wzorów na jej skoki z twierdzenia~\ref{tw:ks-primitive}. Jednoznaczność miar daje równość na wszystkich zbiorach borelowskich. W konsekwencji $\int h\,\mathrm d\mu_F=\int hf\,\mathrm d\mu_g$. Z twierdzenia~\ref{tw:ks-measure} obie strony są wskazanymi całkami KS.
\end{proof}

\begin{twierdzenie}[Oszacowanie błędu zwykłych sum]
\label{tw:ks-quadrature}
Niech $f\in C[a,b]$, $g\in\BV[a,b]$ oraz $P=\{a=x_0<\cdots<x_n=b\}$ będzie dowolnym podziałem o średnicy $h$. Dla dowolnych znaczników $\xi_i\in[x_{i-1},x_i]$ zachodzi
\[
 \left|(\KS)\int_a^b f\,\mathrm dg-
 \sum_{i=1}^n f(\xi_i)\bigl(g(x_i)-g(x_{i-1})\bigr)\right|
 \le\omega_f(h)\,V_a^b(g),
\]
gdzie $\omega_f(h)=\sup\{|f(x)-f(y)|:x,y\in[a,b],\ |x-y|\le h\}$ jest modułem ciągłości. Jeśli $f$ jest $L$-lipschitzowska, błąd jest nie większy niż $LhV_a^b(g)$; jeśli jest h\"olderowska ze stałą $K$ i wykładnikiem $\theta\in(0,1]$, błąd nie przekracza $Kh^\theta V_a^b(g)$.
\end{twierdzenie}
\begin{proof}
Z addytywności całki KS i liniowości mamy
\[
 (\KS)\int_a^b f\,\mathrm dg-S_g(f,\mathcal P)
 =\sum_{i=1}^n(\KS)\int_{x_{i-1}}^{x_i}\bigl(f(t)-f(\xi_i)\bigr)\,\mathrm dg(t).
\]
Funkcja $f-f(\xi_i)$ ma na $[x_{i-1},x_i]$ moduł nieprzekraczający $\omega_f(h)$. Lemat~\ref{lem:ks-var} na każdym przedziale daje oszacowanie przez $\omega_f(h)V_{x_{i-1}}^{x_i}(g)$. Sumując i korzystając z addytywności wahania całkowitego, otrzymujemy żądaną nierówność. Dwa ostatnie oszacowania wynikają z definicji modułu ciągłości.
\end{proof}

\begin{uwaga}[Czego nie daje sama średnica podziału]
Jeżeli $f$ ma skok w punkcie, w którym skacze $g$, mała średnica podziału nie gwarantuje przybliżenia całki przez dowolne sumy. Dla $f=g=\ind{[c,b]}$ sumy na bardzo drobnych podziałach mogą być równe $0$ albo $1$ w zależności od znaczników. Właśnie dlatego w definicji KS stosujemy cechowania zależne od punktu. Oszacowanie z twierdzenia~\ref{tw:ks-quadrature} wymaga ciągłości $f$.
\end{uwaga}

\begin{przyklad}[Przedziałowe oszacowanie dla kwadratury]
Na $[0,1]$ niech $f(t)=\sin t$ i $g$ będzie dowolną niemalejącą funkcją o $g(0)=0$, $g(1)=3$. Dla podziału o średnicy $h=10^{-3}$ dowolna suma Stieltjesa z wartościami $\sin\xi_i$ różni się od całki KS o co najwyżej $3\cdot10^{-3}$, gdyż $|f'|\le1$. Nie potrzebujemy pochodnej $g$ ani wiedzy o położeniu jej skoków.
\end{przyklad}

\subsection*{Granice wzorów i porównanie z całką HK}
\begin{stwierdzenie}[Nie każda całka KS jest całką Lebesgue'a--Stieltjesa]
Dla $g(t)=t$ na $[0,1]$ całka KS jest całką Henstocka--Kurzweila. Istnieją funkcje $f:[0,1]\to\R$, które są KS-całkowalne względem $g$, lecz nie należą do $L^1([0,1],\lambda)$. W szczególności równość z twierdzenia~\ref{tw:ks-measure} nie obowiązuje bez założeń o $f$.
\end{stwierdzenie}
\begin{proof}
Niech $F(0)=0$ i $F(t)=t^2\cos(\pi/t^2)$ dla $t\in(0,1]$ oraz $f(0)=0$, $f(t)=F'(t)=2t\cos(\pi/t^2)+(2\pi/t)\sin(\pi/t^2)$ dla $t>0$. Z twierdzenia Newtona--Leibniza dla całki HK otrzymujemy $(\KS)\int_0^1 f\,\mathrm dt=F(1)-F(0)=-1$. Drugi składnik pochodnej nie jest bezwzględnie całkowalny w otoczeniu zera, ponieważ po podstawieniu $u=\pi/t^2$ jego całka bezwzględna jest porównywalna z $\int^\infty|\sin u|/u\,\mathrm du=\infty$. Tymczasem pierwszy składnik należy do $L^1$. Zatem $\int_0^1|f|\,\mathrm d\lambda=\infty$, więc całka Lebesgue'a z $f$ jako funkcji rzeczywistej nie istnieje.
\end{proof}

\begin{uwaga}[Schemat wyboru wzoru do obliczenia]
Gdy $g$ jest schodkowa, stosujemy twierdzenie~\ref{tw:ks-step}. Gdy $g$ ma wahanie ograniczone, rozkładamy jej miarę według twierdzenia~\ref{tw:ks-decomp}. Gdy $f$ jest absolutnie ciągła, a $g$ jedynie regulowana, najlepiej zastosować wniosek~\ref{wn:ks-ibp-ac}. Przy ciągłej funkcji podcałkowej i całkującej o wahaniu ograniczonym można korzystać z oszacowania kwadraturowego~\ref{tw:ks-quadrature}. Dla dowolnego integratora i funkcji bez powyższej regularności nie istnieje jeden uniwersalny wzór redukujący całkę KS do zwykłej całki Lebesgue'a.
\end{uwaga}

